\documentclass[11pt]{article} 
\usepackage{nopageno}
\setcounter{secnumdepth}{1} % Use 1 for sections, 2 for subsections, etc.
\usepackage[margin=1in]{geometry}
\usepackage{titlesec}
\titleformat{\section}[runin]{\normalfont\bfseries}{\thesection}{1em}{}[]
\titlespacing*{\section}{0pt}{0.1\baselineskip}{1ex}

\titleformat{\subsection}[runin]{\normalfont\bfseries}{\thesubsection}{1em}{}[:]
\titlespacing*{\subsection}{0pt}{0.1\baselineskip}{1ex}

\usepackage{enumitem}
\setlist[itemize]{leftmargin=1.5em, labelsep=0.25em}
\begin{document}

% Remove page numbers---Research.gov will paginate the document instead
%
% \pagestyle{empty} 

\noindent \textbf{Data Management and Sharing Plan} 
Details regarding institutional support for data management are provided in the Facilities, Equipment, and Other Resources document under the ``Research Data Stewardship and Curation" section.

\section*{1. Types of Data, Software, and Curriculum Materials}
This project will generate a multi-layered, AI-ready Student Math Modeling Learning Dataset.
The project anticipates a longitudinal sample of 4,500 high school students (Grades 9-12) primarily from rural Pennsylvania school districts.
The resulting data volume is estimated at tens of millions of annotated events, yielding several tens of gigabytes of raw telemetry. 
Data artifacts fall into one of five main categories:
\vspace{-0.15cm}
\begin{itemize}
    \item \textbf{Process-Level Telemetry:} Captured via an embedded framework in Jupyter notebooks that hooks directly into the IPython kernel. This includes timestamped revision histories, cell-level code diffs, interpreter errors and warnings, execution patterns, and inactivity/idle timers.\vspace{-0.3cm}
    \item \textbf{Assessment and Outcome Data:} Quantitative measures including pre- and post-assessments using validated instruments to track gains in algebra, statistics, and data literacy. \vspace{-0.3cm}
    \item \textbf{Qualitative Artifacts:} Non-code products such as written reflections, variable justifications, model documentation, and final project submissions.\vspace{-0.3cm}
    \item \textbf{Observation and Metadata:} Contextual data including classroom discourse patterns, teacher-student interaction logs, and student demographics. Also includes instructional metadata such as task sequences and standardized Fidelity of Implementation (FoI) ratings used to cross-reference pedagogical variations against student telemetry.\vspace{-0.3cm}
    \item \textbf{Structured Annotations:} Expert-coded labels for reasoning patterns, error types from a statistics education taxonomy, and similar indicators.
\end{itemize}
\vspace{-0.15cm}
Curriculum materials include courseware modules for Algebra and Statistics situated in industrial contexts, e.g., materials science and data center cooling systems.
The project will produce a methodological protocol for porting engineering research data to K-12 curricula.

\section*{2. Standards for Data and Metadata Format}
The project adheres to FAIR (Findable, Accessible, Interoperable, Reusable) principles, storing data in open, non-proprietary formats to ensure compatibility with diverse analysis tools. This includes telemetry and process data (JSON and JSONL), tabular and assessment data (CSV), software and notebooks (standard Python script files and Jupyter notebooks), and configuration and metadata (YAML manifests for the courseware compiler).
Metadata will align with xAPI and Caliper standards where applicable to ensure interoperability with future intelligent tutoring systems.
Metadata will include persistent identifiers (DOIs) for all dataset components and machine-readable descriptions to facilitate discovery through standard search mechanisms.
For domain-specific engineering data used in tasks, the project will leverage existing metadata structures from Penn State's LiST and ScholarSphere.

\section*{3. Policies for Access and Sharing}

In accordance with the NSF Gold Standard Science Implementation Plan, all data will be made available no later than 1 year after the project has finished.
As a justified exception to this unrestricted public sharing requirement to safeguard human subjects, the project utilizes a tiered sharing model to balance analytical richness with student privacy:
\vspace{-0.3cm}
\begin{itemize}
    \item \textbf{Public Tier:} Provides fully de-identified process data and qualitative artifacts. Demographic data in this tier is limited to aggregate categories. This tier is intended for broad research and machine learning applications.\vspace{-0.3cm}
    \item \textbf{Restricted Tier:} Provides fine-grained demographic data necessary for intersectional analysis and longitudinal studies. Access to this tier is granted solely through formal Data Use Agreements (DUAs) and requires evidence of IRB approval from the requesting institution. All data collection and transfer will utilize encrypted protocols. Raw data is stored on Penn State’s secure research infrastructure, with access restricted to IRB-approved personnel.\vspace{-0.3cm}
\end{itemize}
In the context of this collaboration, data will be co-produced and shared among network members (researchers, practitioners, and data scientists) through a secure internal staging environment during the project phase to ensure pedagogical and technical accuracy before public release.

\section*{4. Provisions for Protection of Privacy and Intellectual Property}

The project operates under Penn State IRB approval with informed consent from students and parents/guardians, and all student data sharing and storage will be strictly subject to Family Educational Rights and Privacy Act (FERPA) regulations.
A robust de-identification protocol partitions telemetry into structural and semantic traces. Structural traces (timestamps, error codes, flags) contain no PII by construction.
Semantic traces (markdown cells, written strings) undergo automated scrubbing via pattern-matching algorithms to remove direct identifiers (names, school names, specific dates), followed by manual random sampling for verification.
Written reflections will be redacted before annotation.
Classroom observations will use participant codes; no video or audio recordings are included in shared datasets. 
Intellectual property for software, including the Composable Courseware and telemetry layers, will be managed via GitHub using open-source licenses.
Publicly released data will be licensed under Creative Commons Attribution (CC BY) to maximize reusability while ensuring proper credit to the network.

\section*{5. Policies and Provisions for Re-use, Re-distribution, and Derivatives}

The project encourages the production of derivatives, particularly by other CAMEL networks or educational researchers. To facilitate re-use, the dataset will include: codebooks and annotation rubrics, sample analysis scripts and tutorial notebooks to lower the entry barrier for researchers, explicitly configurable complexity (statistical, structural, and provenance) allowing researchers to tune datasets for different preparation levels.
The modular nature of the Composable Courseware ensures that specific modules can be swapped while maintaining the underlying dataset and model class.

\section*{6. Plans for Archiving and Preservation}

Data will be stored on secure, redundant storage systems maintained by the Pennsylvania State University’s Institute for Computational and Data Sciences (ICDS) Advanced Cyber Infrastructure (ICDS-ACI) during the active project phase. ICDS-ACI provides both active storage (for data that is being worked on, requiring frequent access) and near-line storage (for back-up purposes and data that needs only infrequent access). 
Active storage is achieved through DDN 12KX40 and GS7K flash storage array systems, while near-line storage utilizes Oracle’s FS1 flash storage appliance and a SL8500 Tape Library. Active storage is backed up to the SL8500 daily. ICDS provides long-term archival services using the Oracle SL8500. Multiple copies of archived data are created through Oracle Hierarchical Storage Manager (Oracle HSM) to safeguard against data corruption. All students and faculty will receive training on secure data handling practices. Personally identifiable information will be protected using encryption and access controls. 
Long-term preservation is ensured through the use of established repositories:
\vspace{-0.3cm}
\begin{itemize}
    \item \textbf{The Inter-university Consortium for Political and Social Research (ICPSR)} will host the complete dataset. \vspace{-0.3cm}
    \item \textbf{The Open Science Framework (OSF)} will host supplementary materials, including curriculum tasks, rubrics, and the Data Translation Framework. \vspace{-0.3cm}
    \item \textbf{GitHub} will serve as the active development site for software, with snapshots archived to generate citable DOIs upon project milestones or publication. \vspace{-0.3cm}
\end{itemize}
Peer-reviewed manuscripts will be available via the NSF Public Access Repository (PAR).

PI Napolitano will be responsible for the telemetry framework and dataset conversion.
Co-PI Yao (Science of Learning Lead) will oversee data collection protocols and annotation reliability.
These repositories provide persistent DOIs and ensure data remains accessible and citable beyond the three-year project period.
The project team will coordinate with the Phase II National Coordinator and Collaboratory to facilitate long-term tracking of dataset usage, derivative tools, and associated educational outcomes.
\end{document}